% !TEX program = lualatex
%% Supplementary Materials
\documentclass[5p,twocolumn]{elsarticle}

\usepackage{luatexja}
\usepackage{luatexja-fontspec}
\usepackage{amsmath,amssymb}
\usepackage{graphicx}
\usepackage{booktabs}
\usepackage{multirow}
\usepackage{siunitx}
\usepackage{xcolor}
\usepackage[colorlinks=true,linkcolor=blue,citecolor=blue,urlcolor=blue]{hyperref}
\usepackage{longtable}
\usepackage{array}
\usepackage{float}
\usepackage{enumitem}
\usepackage{subcaption}
\usepackage{listings}
\usepackage{tabularx}
\usepackage{url}

\setlist{nosep,leftmargin=*}

% Japanese fonts
\setmainjfont{IPAexMincho}
\setsansjfont{IPAexGothic}

\lstset{
  basicstyle=\ttfamily\footnotesize,
  breaklines=true,
  frame=single,
  language=SQL,
  keywordstyle=\color{blue}\bfseries,
  commentstyle=\color{green!50!black},
  stringstyle=\color{red!70!black},
  numbers=left,
  numberstyle=\tiny\color{gray},
  tabsize=2,
}

\journal{---}

% Section numbering as S1, S2, ...
\renewcommand{\thesection}{S\arabic{section}}
\renewcommand{\thetable}{S\arabic{table}}
\renewcommand{\thefigure}{S\arabic{figure}}
\renewcommand{\theequation}{S\arabic{equation}}

% =====================================================================
\begin{document}

\begin{frontmatter}

\title{Supplementary Materials for:\\
  Schema-Graph-Constrained Natural Language Interface for\\
  L1$_2$ Intermetallic Compound Exploration}

\author[nims]{Satoshi Minamoto\corref{cor1}}
\cortext[cor1]{Corresponding author}
\ead{minamoto.satoshi@nims.go.jp}
\address[nims]{National Institute for Materials Science (NIMS),
  1-2-1 Sengen, Tsukuba, Ibaraki 305-0047, Japan}

\end{frontmatter}

% =====================================================================
\section{回帰テストカテゴリ詳細}
\label{sec:sup_test_categories}

開発安全性テストとして6カテゴリ39件の回帰テスト、
SQL検証テスト6件、Coverage Score・パーサーテスト15件、
Intent Classifierテスト20件の計80件を構築した（表~\ref{tab:test_categories}）。
全80件がPASSを達成した。

\begin{table*}[htbp]
\centering
\caption{回帰テストカテゴリ（全80件、全PASS）。}
\label{tab:test_categories}
\small
\begin{tabularx}{\textwidth}{llcX}
  \toprule
  カテゴリ & ID & $n$ & 検証焦点 \\
  \midrule
  正常系 & A01--A15 & 15 & L1$_2$元素・プロトタイプ・安定性・ソート・複合条件 \\
  該当なし & B01--B05 & 5 & 希ガス元素、非存在組合せ $\to$ 0行期待 \\
  曖昧な入力 & C01--C10 & 10 & 空入力、無関係テキスト、単語のみ \\
  矛盾条件 & D01--D02 & 2 & 「安定かつ準安定」等 \\
  SQLインジェクション & E01--E05 & 5 & DROP, DELETE, INSERT、ピギーバック攻撃 \\
  安全性検査 & F01--F02 & 2 & 直接SQL入力（SELECT, UPDATE） \\
  SQL検証 & --- & 6 & カラム/テーブルホワイトリスト、JOIN検証 \\
  Coverage・パーサー & --- & 15 & 数値条件、化学式、Coverage Score \\
  Intent Classifier & --- & 20 & VASP/DFTワークフロー検出、分類 \\
  \bottomrule
\end{tabularx}
\end{table*}

% =====================================================================
\section{回帰テストケース一覧}
\label{sec:sup_all_tests}

表~\ref{tab:all_tests}に回帰テスト44件（正常系・該当なし・曖昧・矛盾・
SQLインジェクション・安全性）のクエリと結果を示す。
全評価クエリ100件（Easy 20、Medium 30、Hard 30、Very Hard 20）の
gold SQL、期待結果セット、5手法比較結果、
プロンプトテンプレート、
\texttt{allowed\_schema.yaml}、\texttt{material\_terms.yaml}は
検証資材リポジトリに収録されている。

\begin{table*}[htbp]
\centering
\caption{回帰テストケース一覧（全44件、全PASS）。}
\label{tab:all_tests}
\footnotesize
\begin{tabular}{llp{8cm}l}
  \toprule
  ID & カテゴリ & 自然言語クエリ & 結果 \\
  \midrule
  A01 & 正常系 & Feを含むL1$_2$化合物を出して & PASS \\
  A02 & 正常系 & 安定なL1$_2$化合物を形成エネルギー順に出して & PASS \\
  A03 & 正常系 & NiとAlの両方を含む化合物を出して & PASS \\
  A04 & 正常系 & L1$_2$化合物の全リストを出して & PASS \\
  A05 & 正常系 & 準安定なL1$_2$化合物を出して & PASS \\
  A06 & 正常系 & Coを含む安定なL1$_2$化合物を出して & PASS \\
  A07 & 正常系 & FeとAlを含むL1$_2$化合物を出して & PASS \\
  A08 & 正常系 & $\gamma'$化合物を出して & PASS \\
  A09 & 正常系 & ニッケルを含むL1$_2$型化合物を出して & PASS \\
  A10 & 正常系 & L1$_2$化合物の全データを出して & PASS \\
  A11 & 正常系 & Tiを含むL1$_2$化合物を格子定数順（降順）に出して & PASS \\
  A12 & 正常系 & 安定なL1$_2$で格子定数3.5\AA 以上のものを出して & PASS \\
  A13 & 正常系 & Cu$_3$Au型化合物を出して & PASS \\
  A14 & 正常系 & 鉄を含むL1$_2$化合物を出して & PASS \\
  A15 & 正常系 & 安定なL1$_2$化合物でNiを含むものを出して & PASS \\
  \midrule
  B01 & 該当なし & Xeを含むL1$_2$化合物を出して & PASS \\
  B02 & 該当なし & UとPuを含むL1$_2$化合物を出して & PASS \\
  B03 & 該当なし & RnとOgを含むL1$_2$化合物を出して & PASS \\
  B04 & 該当なし & ScとIrを含む安定なL1$_2$化合物を出して & PASS \\
  B05 & 該当なし & PtとGaを含む安定なL1$_2$化合物を出して & PASS \\
  \midrule
  C01 & 曖昧 & 安定な化合物を出して & PASS \\
  C02 & 曖昧 & L12 & PASS \\
  C03 & 曖昧 & 安定なものを出して & PASS \\
  C04 & 曖昧 & （空入力） & PASS \\
  C05 & 曖昧 & 今日の天気を教えて & PASS \\
  C06 & 曖昧 & Fe & PASS \\
  C07 & 曖昧 & 格子定数 & PASS \\
  C08 & 曖昧 & 形成エネルギーが低い安定なL1$_2$化合物 & PASS \\
  C09 & 曖昧 & NiAl L12 & PASS \\
  C10 & 曖昧 & NiのL12 & PASS \\
  \midrule
  D01 & 矛盾 & 安定かつ準安定なL1$_2$化合物 & PASS \\
  D02 & 矛盾 & NiなしのNi L1$_2$化合物 & PASS \\
  \midrule
  E01 & 拒否 & DROP TABLE material\_entry; & PASS \\
  E02 & 拒否 & SELECT *; DELETE FROM composition; & PASS \\
  E03 & 拒否 & L1$_2$化合物; DROP TABLE structure; & PASS \\
  E04 & 拒否 & SELECT * FROM secret\_passwords & PASS \\
  E05 & 拒否 & INSERT INTO material\_entry VALUES (...) & PASS \\
  \midrule
  F01 & 安全性 & SELECT entry\_id FROM material\_entry & PASS \\
  F02 & 安全性 & UPDATE material\_entry SET formula='X' & PASS \\
  \bottomrule
\end{tabular}
\end{table*}

% =====================================================================
\section{生成SQL例}
\label{sec:sup_sql_examples}

代表的なクエリに対する各手法の生成SQL例を示す。

\lstset{basicstyle=\ttfamily\scriptsize,breaklines=true,columns=flexible,numbers=none,frame=single,xleftmargin=0pt,xrightmargin=0pt,linewidth=\columnwidth}

\subsection{「Niを含む安定なL1$_2$化合物を出して」}

\paragraph{LLM-only（B1）：}
\begin{lstlisting}
SELECT formula, stability
FROM materials
WHERE element = 'Ni'
  AND structure_type = 'L12'
  AND stability = 'stable';
-- テーブル幻覚: materials, stability等は
-- 存在しないテーブル -> 実行不可
\end{lstlisting}

\paragraph{Proposed手法（P）：}
\begin{lstlisting}
SELECT DISTINCT m.entry_id, m.formula,
  s.prototype, s.lattice_a, s.space_group
FROM material_entry m
  JOIN composition c
    ON c.entry_id = m.entry_id
  JOIN structure s
    ON s.entry_id = m.entry_id
  JOIN phase_stability ps
    ON ps.entry_id = m.entry_id
WHERE (s.prototype = 'L12'
    OR s.strukturbericht = 'L12')
  AND c.element = 'Ni'
  AND ps.energy_above_hull = 0
LIMIT 10000;
-- Schema Graph走査で必要な3 JOINのみ
-- テーブル幻覚0%, DISTINCT, LIMIT付き
\end{lstlisting}

\subsection{「NiとAlの両方を含むL1$_2$化合物を出して」}

\paragraph{Full Schema（B2）：}
\begin{lstlisting}
SELECT m.entry_id, m.formula
FROM material_entry m
  JOIN composition c
    ON c.entry_id = m.entry_id
  JOIN structure s
    ON s.entry_id = m.entry_id
WHERE c.element = 'Ni'
  AND c.element = 'Al'
  AND s.prototype = 'L12';
-- 致命的: 単一行が両条件を同時に
-- 満たすことは不可能 -> 常に0行
\end{lstlisting}

\paragraph{Proposed手法（P）：}
\begin{lstlisting}
SELECT DISTINCT m.entry_id, m.formula
FROM material_entry m
  JOIN structure s
    ON s.entry_id = m.entry_id
WHERE (s.prototype = 'L12'
    OR s.strukturbericht = 'L12')
AND EXISTS (
  SELECT 1 FROM composition
  WHERE entry_id = m.entry_id
    AND element = 'Ni')
AND EXISTS (
  SELECT 1 FROM composition
  WHERE entry_id = m.entry_id
    AND element = 'Al')
LIMIT 10000;
-- EXISTS副問い合わせで正確な多元素AND
-- -> Ni3Al (1行) を正しく返却
\end{lstlisting}

% =====================================================================
\section{SQLインジェクション・安全性テスト結果}
\label{sec:sup_injection}

7件の敵対的入力（E01--E05、F01--F02）はすべてSQLガードにより遮断された
（表~\ref{tab:injection_results}）。

\begin{table}[htbp]
  \centering
  \caption{SQLインジェクション・安全性テスト結果（全件遮断）。}
  \label{tab:injection_results}
  \small
  \begin{tabularx}{\linewidth}{lXl}
    \toprule
    ID & 入力 & 遮断理由 \\
    \midrule
    E01 & \texttt{DROP TABLE material\_entry;} & 禁止キーワード：DROP \\
    E02 & \texttt{SELECT *; DELETE FROM composition;} & 複数文 + DELETE \\
    E03 & L1$_2$化合物; \texttt{DROP TABLE structure;} & 複数文 + DROP \\
    E04 & \texttt{SELECT * FROM secret\_passwords} & 非許可テーブル \\
    E05 & \texttt{INSERT INTO material\_entry ...} & 禁止キーワード：INSERT \\
    F01 & \texttt{SELECT entry\_id FROM material\_entry} & LIMITなし $\to$ 自動付加 \\
    F02 & \texttt{UPDATE material\_entry SET ...} & 禁止キーワード：UPDATE \\
    \bottomrule
  \end{tabularx}
\end{table}

% =====================================================================
\section{Ni$_3$Al近傍格子定数候補}
\label{sec:sup_lattice_match}

$|a - a_{\text{Ni}_3\text{Al}}| \leq 0.03$~\AA を満たす候補を
表~\ref{tab:lattice_match}に示す。

\begin{table}[htbp]
  \centering
  \caption{Ni$_3$Al近傍格子定数候補（$|\Delta a| \leq 0.03$~\AA）。}
  \label{tab:lattice_match}
  \small
  \begin{tabular}{lcccc}
    \toprule
    化合物 & $a$ (\AA) & $|\Delta a|$ & $E_\text{hull}$ & $E_f$ \\
    \midrule
    Ni$_3$Al & 3.572 & 0.002 & 0.000 & $-0.420$ \\
    Ni$_3$Zn & 3.560 & 0.010 & 0.032 & $-0.340$ \\
    Pt$_3$Sn & 3.560 & 0.010 & 0.032 & $-0.490$ \\
    Cu$_3$Nb & 3.560 & 0.010 & 0.004 & $-0.315$ \\
    Ni$_3$V & 3.580 & 0.010 & 0.036 & $-0.345$ \\
    Cu$_3$Ta & 3.580 & 0.010 & 0.008 & $-0.320$ \\
    Pt$_3$Zn & 3.580 & 0.010 & 0.036 & $-0.495$ \\
    Co$_3$Ti & 3.550 & 0.020 & 0.000 & $-0.450$ \\
    Cu$_3$Ti & 3.540 & 0.030 & 0.000 & $-0.310$ \\
    Pt$_3$Si & 3.540 & 0.030 & 0.028 & $-0.485$ \\
    \bottomrule
  \end{tabular}
\end{table}

格子整合性に優れた候補の中で、
Cu$_3$Nb（$E_\text{hull} = 0.004$）とCu$_3$Ta（$E_\text{hull} = 0.008$）は
準安定であり、Co$_3$Ti（$E_\text{hull} = 0.0$）とCu$_3$Ti（$E_\text{hull} = 0.0$）は
安定相として確認された。
これらはNi$_3$Alとの格子ミスマッチが小さく、
析出強化の観点から有望な候補である。

% =====================================================================
\section{安定L1$_2$候補の抽出}
\label{sec:sup_stable_l12}

$E_\text{hull} \leq 0.05$~eV/atomを満たす安定・準安定L1$_2$候補を抽出した。
安定（$E_\text{hull} = 0$）12件、準安定（$0 < E_\text{hull} \leq 0.05$）45件の
計57件が候補として抽出された。
安定相のうち、最も負の形成エネルギーを持つのは
Pt$_3$Al（$E_f = -0.55$~eV/atom）であり、
次いでAl$_3$Sc（$-0.50$）、Co$_3$Ti（$-0.45$）、Pt$_3$Ge（$-0.45$）が続く。

% =====================================================================
\section{LLMモード設定}
\label{sec:sup_llm_config}

LLMモード（Level 2）ではgpt-5.5を使用した。
構造化プロンプトは以下の4要素から構成される：

\begin{enumerate}
  \item \textbf{システムメッセージ}：タスク説明 + スキーマYAML（テーブル名、カラム名、型、FK関係）
  \item \textbf{Few-Shot事例}：RAGストアからのTop-$k$類似NL--SQLペア
  \item \textbf{ユーザーメッセージ}：自然言語クエリ
  \item \textbf{制約}：「SELECT文のみ生成」「LIMIT 10000を含める」
        「スキーマに記載されたテーブルのみ使用」
\end{enumerate}

プロンプトテンプレート（\texttt{llm/prompt\_templates/sql\_generation\_prompt.md}）
および\texttt{allowed\_schema.yaml}は検証資材リポジトリに収録されている。

\end{document}
